% BEGIN REV09_FORMA_PARTIDA_ENERGIA
\section{Forme déployée, intégralité et énergie réticulaire}
\label{corpus9:xii:forma-partida}

La déformation du théorème~\ref{rec:thm:red} part du réseau marqué
\(\Lambda\), de la direction \(u=P_3P_{11}K\) et des coordonnées
régionales préalablement engendrées. Nous notons \(E\) son espace réel
de dimension vingt-quatre et \(\langle\,\cdot\,,\,\cdot\,\rangle_B\)
la métrique positive des douze plans orthogonaux de type \(A_2\).
Nous conservons la normalisation de \eqref{rec:eq:flujo} et
\eqref{rec:eq:red} :
\begin{equation}
 \kappa=\frac{\log\rho}{\|u\|},\qquad h_i=\kappa u_i,
 \qquad
 G_t=\bigoplus_{i=1}^{12}e^{th_i}I_{A_{2,i}},\qquad
 \Lambda_t=G_t\Lambda.
 \label{corpus9:xii:normalizacion-flujo}
\end{equation}
L’opérateur \(G_t\) est autoadjoint pour \(B\), et
\(G_t^{-1}=G_{-t}\). Si l’on écrit la même collection avec le paramètre
non normalisé, \(D_s=\bigoplus_i e^{su_i}I_{A_{2,i}}\), la relation
exacte est \(G_t=D_{\kappa t}\). Le paramètre de déformation et
l’horloge TPK maintiennent leurs domaines respectifs.

\subsection{Variation arithmétique de la projection euclidienne}
\label{corpus9:xii:variacion-euclidea}

Le système de coordonnées ordonné du lemme~\ref{lem:k-sector-no-nulo} détermine,
par la somme finie \eqref{eq:k-proyectores-incidencia},
\begin{equation}
 u_1=-\frac{495}{4}-\frac{233}{10}\sqrt5,
 \qquad \sum_{i=1}^{12}u_i=0.
 \label{corpus9:xii:coordenada-testigo}
\end{equation}
La normalisation régionale satisfait \(\rho=\pi_{\HMT}/\varphi_{\HMT}^2>1\) :
dans la réalisation déjà établie, \(\pi_{\HMT}>3\) et
\(\varphi_{\HMT}^2=(3+\sqrt5)/2<3\). Par conséquent, \(\kappa>0\).
Ces comparaisons sont effectuées sur les sorties régionales, après
leur génération.

\begin{proposition}[Témoin local de variation de l’intégralité
]
\label{corpus9:xii:testigo-integralidad}
Il existe \(\delta>0\) tel que, pour \(0<|t|<\delta\), le réseau
\(\Lambda_t\) contient un vecteur de norme au carré non entière.
En conséquence, sa métrique euclidienne ne possède ni l’intégralité ni une
isométrie avec le réseau de Leech.
\end{proposition}
\begin{proof}
Le vecteur marqué de \eqref{exc:vector} est
\(v_\alpha=4\rho_1+\sum_{j=2}^{12}\rho_j\), avec
\(\|\rho_j\|_B^2=2\) dans des plans mutuellement orthogonaux et
\(v_\alpha^2=54\). La construction du voisin
\eqref{exc:vecino} donne \(v_\alpha/3\in\Lambda\). Sa norme transportée est

\begin{equation}
 f(t)=\left\|G_t\frac{v_\alpha}{3}\right\|_B^2
 =\frac29\left(16e^{2th_1}+\sum_{j=2}^{12}e^{2th_j}\right).
 \label{corpus9:xii:norma-testigo}
\end{equation}
La somme nulle et \eqref{corpus9:xii:coordenada-testigo} produisent
\begin{equation}
 \begin{aligned}
 f(0)&=6,\\
 f'(0)&=\frac49\left(16h_1+\sum_{j=2}^{12}h_j\right)
       =\frac{20}{3}\kappa u_1
       =\kappa\left(-825-\frac{466}{3}\sqrt5\right)\ne0.
 \end{aligned}
 \label{corpus9:xii:derivada-testigo}
\end{equation}
La continuité de \(f'\) fournit un intervalle
\((-\delta,\delta)\) dans lequel \(f\) est strictement monotone.
En réduisant \(\delta\), on obtient en outre
\(11/2<f(t)<13/2\). Pour \(0<|t|<\delta\), la monotonie donne
\(f(t)\ne f(0)=6\) ; le seul entier dans cet intervalle est six.
Ainsi, \(G_t(v_\alpha/3)\) a une norme au carré non entière. L’intégralité
d’un réseau exige des normes au carré entières, et les isométries
conservent ces normes. L’intégralité de Leech prouve la dernière
conclusion.
\end{proof}

La déformation conserve le covolume et la symétrie ternaire tandis que
sa structure métrique arithmétique varie. La réciprocité
\(\Lambda_t^*=\Lambda_{-t}\) utilise l’autoadjonction du flot et
l’autodualité initiale. Une permutation ambiante qui échange
dilatations et contractions agit sur le réseau recollé lorsqu’elle
préserve en outre le code et le voisin qui le construisent.


\subsection{Autodualité entière pour la forme de signature déployée}
\label{corpus9:xii:autodualidad-partida}

La considération simultanée des deux projections utilise l’espace
\(E\oplus E\), muni de la forme bilinéaire
\begin{equation}
 \mathcal B((x,p),(x',p'))
 =\langle x,p'\rangle_B+\langle p,x'\rangle_B.
 \label{corpus9:xii:forma-bilineal-partida}
\end{equation}
On définit
\begin{equation}
 \begin{aligned}
 \mathcal G_t&=G_t\oplus G_{-t},\\
 \Gamma_t&=\mathcal G_t(\Lambda\oplus\Lambda)
 =\{(G_t\lambda,G_{-t}\mu):\lambda,\mu\in\Lambda\}.
 \end{aligned}
 \label{corpus9:xii:red-partida}
\end{equation}

\begin{theorem}[Intégralité, parité et autodualité du réseau déployé]
\label{corpus9:xii:teorema-partido}
La forme \(\mathcal B\) a pour signature \((24,24)\).
Pour tout \(t\in\mathbb R\), \(\mathcal G_t\) est une isométrie
de \(\mathcal B\), et \(\Gamma_t\) est un réseau entier, pair et
autodual relativement à cette forme.
\end{theorem}
\begin{proof}
Les coordonnées
\[
 p_L=\frac{x+p}{\sqrt2},\qquad p_R=\frac{x-p}{\sqrt2}
\]
constituent un changement linéaire inversible, avec
\(x=(p_L+p_R)/\sqrt2\) et \(p=(p_L-p_R)/\sqrt2\). Ils satisfont

\[
 \mathcal B((x,p),(x,p))=\|p_L\|_B^2-\|p_R\|_B^2.
\]
La positivité de \(B\) en dimension vingt-quatre donne la signature
\((24,24)\). L’autoadjonction de \(G_t\) et l’identité
\(G_tG_{-t}=I\) impliquent

\[
 \langle G_tx,G_{-t}p'\rangle_B
 =\langle x,G_tG_{-t}p'\rangle_B=\langle x,p'\rangle_B.
\]
L’autre produit croisé satisfait la même égalité. Par conséquent,
\(\mathcal B(\mathcal G_t\xi,\mathcal G_t\eta)
=\mathcal B(\xi,\eta)\).

Dans \(\Gamma_0=\Lambda\oplus\Lambda\), les produits croisés sont
entiers, car \(\Lambda\) est entier, et
\[
 \mathcal B((\lambda,\mu),(\lambda,\mu))
 =2\langle\lambda,\mu\rangle_B\in2\mathbb Z.
\]
Cela prouve l’intégralité et la parité. Si \((x,p)\) appartient au dual
de \(\Gamma_0\) par rapport à \(\mathcal B\), ses appariements
avec \((\lambda,0)\) et \((0,\mu)\), pour tout
\(\lambda,\mu\in\Lambda\), exigent respectivement
\(p,x\in\Lambda^*=\Lambda\). Réciproquement, ces deux appartenances
garantissent tous les appariements entiers. Ainsi,
\(\Gamma_0^{*\mathcal B}=\Gamma_0\). L’isométrie
\(\mathcal G_t\) transporte les trois propriétés et donne
\(\Gamma_t^{*\mathcal B}=\mathcal G_t\Gamma_0^{*\mathcal B}
=\Gamma_t\).
\end{proof}

La métrique détermine le domaine de chaque conclusion. Les produits
croisés de \(\mathcal B\) conservent l’intégralité de
\(\Gamma_t\), même pour les paramètres de la
proposition~\ref{corpus9:xii:testigo-integralidad}. Les réseaux
\(\Gamma_t\), considérés avec cette forme déployée, sont isométriques
entre eux. Leurs projections euclidiennes \(\Lambda_t\) et
\(\Lambda_{-t}\), munies de \(B\), varient avec le paramètre.

\begin{proposition}[Échange des projections réciproques]
\label{corpus9:xii:involucion-partida}
L'involution \(\mathcal J(x,p)=(p,x)\) satisfait
\[
 \mathcal J\Gamma_t=\Gamma_{-t},\qquad
 \mathcal J\mathcal G_t\mathcal J=\mathcal G_{-t}.
\]
Dans les coordonnées \((p_L,p_R)\), il fixe \(p_L\) et inverse
le signe de \(p_R\).

\end{proposition}
\begin{proof}
L’échange transforme \((G_t\lambda,G_{-t}\mu)\) en
\((G_{-t}\mu,G_t\lambda)\). Lorsque l’on parcourt \(\lambda,\mu\in\Lambda\),
l’image parcourt \(\Gamma_{-t}\). La conjugaison échange les
deux blocs de \(\mathcal G_t\), produisant \(\mathcal G_{-t}\).
Enfin, \(x+p\) est conservé et \(x-p\) change de signe.

\end{proof}

Cette involution agit sur la construction réticulaire complète. Une
représentation physique ultérieure doit aussi transporter les opérateurs
et les données d’identification propres à cette réalisation.

\subsection{Opérateur positif et équivalence unitaire des domaines
}
\label{corpus9:xii:energia-reticular}

L’énergie des deux projections est définie sur le réseau fixe par
\begin{equation}
 \mathsf E_t^{\rm lat}(\lambda,\mu)
 =\|G_t\lambda\|_B^2+\|G_{-t}\mu\|_B^2,
 \qquad \lambda,\mu\in\Lambda.
 \label{corpus9:xii:energia}
\end{equation}
Elle est positive ou nulle, s’annule exactement en \((0,0)\) et satisfait
\(\mathsf E_t^{\rm lat}(\lambda,\mu)
=\mathsf E_{-t}^{\rm lat}(\mu,\lambda)\). L’échange des deux
contributions exprime, dans les douze plans marqués, la réciprocité
d’une somme de termes d’échelles opposées. Pour
\(M_h=\max_i|h_i|\), la décomposition orthogonale en plans donne

\begin{equation}
 e^{-2|t|M_h}\|v\|_B^2\leq\|G_tv\|_B^2
                    \leq e^{2|t|M_h}\|v\|_B^2.
 \label{corpus9:xii:cotas-cuadraticas}
\end{equation}
En effet, chaque terme \(\|v_i\|_B^2\) est multiplié par
\(e^{2th_i}\), compris entre ces deux bornes.

\begin{theorem}[Caractère autoadjoint, résolvante compacte et dualité unitaire]
\label{corpus9:xii:operador-positivo}
Dans \(\ell^2(\Lambda\oplus\Lambda)\), l’opérateur
\[
 (H_t^{\rm lat}\psi)(\lambda,\mu)
 =\mathsf E_t^{\rm lat}(\lambda,\mu)\psi(\lambda,\mu)
\]
de domaine maximal
\begin{equation}
 \mathcal D(H_t^{\rm lat})=
 \left\{\psi\in\ell^2(\Lambda\oplus\Lambda):
 \sum_{\lambda,\mu}\bigl(\mathsf E_t^{\rm lat}(\lambda,\mu)\bigr)^2
                  |\psi(\lambda,\mu)|^2<\infty\right\}
 \label{corpus9:xii:dominio-operador}
\end{equation}
est positif, auto-adjoint et à résolvante compacte. L’application
\(U\psi(\lambda,\mu)=\psi(\mu,\lambda)\) est unitaire et vérifie
\begin{equation}
 U\mathcal D(H_t^{\rm lat})=\mathcal D(H_{-t}^{\rm lat}),\qquad
 UH_t^{\rm lat}U^{-1}=H_{-t}^{\rm lat}.
 \label{corpus9:xii:equivalencia-unitaria}
\end{equation}
\end{theorem}
\begin{proof}
Les fonctions à support fini appartiennent au domaine et forment un
sous-espace dense. L’énergie est réelle. Si \(\psi\) appartient au domaine
de l’adjoint, évaluer sa fonctionnelle contre chaque vecteur de la base
orthonormale à support unitaire impose que la coordonnée de l’adjoint
soit \(\mathsf E_t^{\rm lat}(\lambda,\mu)\psi(\lambda,\mu)\).
L’appartenance de cette suite à \(\ell^2\) est exactement la condition
\eqref{corpus9:xii:dominio-operador}. Réciproquement, cette condition
permet l’appariement avec toute fonction du domaine par
Cauchy--Schwarz. Le domaine de l’adjoint et son action coïncident avec ceux
de \(H_t^{\rm lat}\), ce qui prouve le caractère autoadjoint. En outre,
\[
 \langle\psi,H_t^{\rm lat}\psi\rangle
 =\sum_{\lambda,\mu}\mathsf E_t^{\rm lat}(\lambda,\mu)
                         |\psi(\lambda,\mu)|^2\geq0.
\]
Par \eqref{corpus9:xii:cotas-cuadraticas},
\[
 \mathsf E_t^{\rm lat}(\lambda,\mu)
 \geq e^{-2|t|M_h}
       \bigl(\|\lambda\|_B^2+\|\mu\|_B^2\bigr).
\]
Le caractère discret de \(\Lambda\oplus\Lambda\) implique que chaque sous-niveau
de l’énergie est fini. La résolvante \((H_t^{\rm lat}+I)^{-1}\)
est diagonale, de coefficients \((1+\mathsf E_t^{\rm lat})^{-1}\).
Sa restriction au sous-niveau \(\mathsf E_t^{\rm lat}\leq R\) est de
rang fini et le reste est de norme au plus \((1+R)^{-1}\).
La résolvante est donc une limite en norme d’opérateurs de rang
fini et est compacte.

L’échange d’indices conserve la norme de \(\ell^2\) et satisfait
\(U^{-1}=U\). En remplaçant \((\lambda,\mu)\) par \((\mu,\lambda)\)
dans la somme qui définit le domaine et en utilisant l’égalité des énergies,
on obtient la première identité de
\eqref{corpus9:xii:equivalencia-unitaria}. Sur ce domaine,

\[
 (UH_t^{\rm lat}U^{-1}\psi)(\lambda,\mu)
 =\mathsf E_t^{\rm lat}(\mu,\lambda)\psi(\lambda,\mu)
 =\mathsf E_{-t}^{\rm lat}(\lambda,\mu)\psi(\lambda,\mu),
\]
ce qui prouve l’identité des opérateurs avec leurs domaines.

\end{proof}

Les trois conservations sont ainsi typées. La métrique euclidienne
conserve le covolume et la symétrie ternaire et relie les paramètres
opposés au moyen du réseau dual. La forme déployée conserve l’intégralité,
la parité et l’autodualité du réseau dupliqué. La réalisation opératorielle
conserve l’énergie au moyen d’une équivalence unitaire qui transporte
son domaine. La lecture thêta du théorème~\ref{rec:thm:theta} exprime
la réciprocité dans la somme convergente des longueurs, et le prolongement
de Mellin conserve sa structure polaire par le
théorème~\ref{rec:thm:mellin}.

% END REV09_FORMA_PARTIDA_ENERGIA
